\documentclass[12pt, a4paper]{article}
\usepackage[utf8]{inputenc}
\usepackage[lithuanian]{babel}
\usepackage[T1]{fontenc}
\usepackage{amsmath, amssymb}
\usepackage{geometry}
\usepackage{enumitem}

\geometry{top=2cm, bottom=2cm, left=2cm, right=2cm}

\title{\textbf{Užduočių mixas 1}}

\date{}

\begin{document}

\maketitle

\section{Sąlygos}

\begin{enumerate}[label=\textbf{\arabic* užduotis.}, leftmargin=*]

    \item Išspręskite lygtį:
    \[
    2^{x+2}\cdot 4^{x+4}\cdot 8^{x+8}
    =\left(\frac{1}{16}\right)^{-3}.
    \]

    \item Išskirkite pilną kvadratą:
    \begin{enumerate}[label=\alph*)]
        \item $x^2-12x+10$;
        \item $4x^2y^2-9xy+1$;
        \item $2x^2-\dfrac14 x+5$;
        \item $-3x^2+x-9$.
    \end{enumerate}

    \item Apskaičiuokite nenaudodami smegenų džiovintuvo:
    \[
    \frac{(666^2-34^2)(332\cdot 32+90000)}
    {(10\sqrt7)^2(332^3+300^3)}
    :\frac{1}{\pi}.
    \]
    
    \item Išskaidykite dauginamaisiais:
\begin{enumerate}[label=\alph*)]
    \item $\displaystyle \frac{8}{125}x ^{3} - \frac{1}{64}  $;
    \item $\displaystyle 0,001y ^{3} + 2 \frac{10}{27}   $.

\end{enumerate}

    \item Išspręskite lygtį:
    \[
    \frac{x^3-8}
    {(x^2-4)\left(x+\dfrac4x+2\right)x}
    =\frac34.
    \]

    \item Suprastinkite:
    \[
    \sqrt[3]{\frac{(2x+1)^3-12(x+3)(x-4)-5(3{,}6x+29)}{x^3}}.
    \]

    \item \textbf{[Desertas]} Apskaičiuokite:
    \[
    \frac{1}{1\cdot2}+\frac{1}{2\cdot3}+\cdots+\frac{1}{2025\cdot2026}.
    \]

\end{enumerate}

\newpage

\section{Atsakymai}

\noindent\textbf{1 užduotis.} $x=-\frac{11}{3}$.\par

\noindent\textbf{2 užduotis.} a) $(x-6)^2-26$;\quad b) $\left(2xy-\frac94\right)^2-\frac{65}{16}$;\quad c) $2\left(x-\frac1{16}\right)^2+\frac{639}{128}$;\quad d) $-3\left(x-\frac16\right)^2-\frac{107}{12}$.\par

\noindent\textbf{3 užduotis.} $\pi$.\par

\noindent\textbf{4 užduotis.} a) $\left(\frac{2}{5}x-\frac14\right)\left(\frac{4}{25}x^2+\frac{1}{10}x+\frac{1}{16}\right)$;\quad b) $\left(\frac{1}{10}y+\frac43\right)\left(\frac{1}{100}y^2-\frac{2}{15}y+\frac{16}{9}\right)$.\par

\noindent\textbf{5 užduotis.} $x=6$.

\noindent\textbf{6 užduotis.} $2$, kai $x\ne0$.

\noindent\textbf{7 užduotis.} $\frac{2025}{2026}$.

\end{document}
